\chapter{Referências e bibliografia}
\label{chap:referencias}

\section{Separando dados bibliográficos}

Referências\index{bibliografia} não devem ser digitadas e numeradas manualmente. Um arquivo
\arquivo{.bib} armazena cada fonte como um registro identificado por uma chave:

\begin{lstlisting}
@book{lamport1994,
  author    = {Leslie Lamport},
  title     = {LaTeX: A Document Preparation System},
  year      = {1994},
  publisher = {Addison-Wesley}
}
\end{lstlisting}

No texto, \comando{citep} produz uma citação entre parênteses. O comando
\comando{citet} incorpora o autor à frase quando usamos \arquivo{natbib}. Ao
final de \arquivo{main.tex}, indicamos o estilo e o banco:

\begin{lstlisting}[language=LaTeXBook]
\bibliographystyle{plainnat}
\bibliography{bibliography}
\end{lstlisting}

Somente fontes citadas aparecem na lista. Use chaves curtas e estáveis, seguindo
o padrão \emph{sobrenome, ano e tema}. Complete autor, título, ano e publicação;
para páginas web, registre endereço e data de acesso.

\section{O ciclo de compilação}

A bibliografia exige mais de uma etapa: LaTeX descobre as citações, BibTeX monta
a lista e LaTeX incorpora e numera o resultado. \arquivo{latexmk} detecta essa
sequência. Se aparecer \verb|[?]|, compile novamente e procure no log por uma
chave inexistente ou erro no arquivo \arquivo{.bib}.

Projetos acadêmicos podem preferir \arquivo{biblatex} com Biber, que oferece
controle mais amplo. Para começar, \arquivo{natbib} e BibTeX formam um fluxo
simples e suficiente. Não misture os dois sistemas na mesma configuração.

\section{Índice remissivo}

O índice remissivo ajuda o leitor a localizar conceitos, não apenas títulos.
Carregue \arquivo{makeidx} e ative a criação do índice no preâmbulo. Marque os
termos no texto e imprima a lista no final do documento.

\begin{lstlisting}[language=LaTeXBook]
LuaLaTeX\index{LuaLaTeX} permite usar fontes OpenType.
\end{lstlisting}

Marque apenas ocorrências úteis. Indexar toda repetição cria ruído. Assim como a
bibliografia, o índice depende de processamento auxiliar que o \arquivo{latexmk}
normalmente executa.

\begin{pratica}
  Adicione um livro e uma página web ao seu \arquivo{bibliography.bib}. Cite-os
  em um parágrafo e marque três conceitos para o índice remissivo. Compile e
  confirme as duas listas no final do PDF.
\end{pratica}
